\chapter{Evaluation and Reflection}
\label{ch:evaluation}

\section{Designing an evaluation when there is no ground truth}
\label{sec:eval-design}

The obvious evaluation is accuracy against real outcomes, and mostly it is not
available. Fusion21's Doc 4 states that most procurement disputes settle before
proceedings, so cases producing a published judgment are already a biased
minority; of the 23 here, nine were decided on an interim application where the
merits were never reached. Any single accuracy number over that set would be both
small and misleading.

The evaluation is built instead as a stack of partial checks
(Figure~\ref{fig:evidence}), each answering a question the others can't, and each
reported with the sample size that produced it, following the reporting
discipline Dodge et al.\ \cite{dodge2019showyourwork} argue for. Every figure below comes from an
analysis script in the repository run over committed logs, and every script was
re-run immediately before submission. At that point the corpus stood at 299
completed negotiations across 59 batches.

\begin{figure}[htbp]
\centering
\resizebox{\textwidth}{!}{%
\begin{tikzpicture}[font=\footnotesize,
 band/.style={draw=f21ink,rounded corners=2pt,line width=0.7pt,align=left,
 inner sep=5pt,text width=11.6cm}]
 \node[band,fill=f21mist] (a)
 {\textbf{Does it produce well-formed output at all?}\quad
 Structural compliance over 2{,}865 structured responses~(\S\ref{sec:structural})};
 \node[band,below=2.6mm of a,fill=f21mist] (b)
 {\textbf{Does it agree with reality where reality is knowable?}\quad
 14 merits dispositions~(\S\ref{sec:agreement})};
 \node[band,below=2.6mm of b,fill=f21green!10] (c)
 {\textbf{Does the architecture earn its complexity?}\quad
 Three baselines, one ablation, $n=6$ and $n=18$~(\S\ref{sec:baselines})};
 \node[band,below=2.6mm of c,fill=f21green!10] (d)
 {\textbf{Does a prompt change alter judgement or only confidence?}\quad
 Controlled V3/V4 ablation, $n=8$ per cell, Fisher's exact~(\S\ref{sec:ablation})};
 \node[band,below=2.6mm of d,fill=f21amber!10] (e)
 {\textbf{Is the reasoning grounded in what the model was given?}\quad
 Three independent fabrication screens~(\S\ref{sec:fabrication})};
 \node[band,below=2.6mm of e,fill=f21amber!10] (f)
 {\textbf{Is the output usable by the person who has to act on it?}\quad
 Concession dynamics, BATNA self-report, readability~(\S\ref{sec:dynamics}--\S\ref{sec:readability})};
\end{tikzpicture}}
\caption[Evaluation design]{The evaluation as a stack of partial checks. No single
row is sufficient; the fabrication screens in particular exist because a system can
score well on every row above them while inventing the evidence it reasons from.}
\label{fig:evidence}
\end{figure}

\section{The corpus}
\label{sec:corpus}

Twenty-three real UK procurement judgments were researched, individually verified
against BAILII or the National Archives case law service, and written into the
project as dense source texts, then run through the same extraction path a user
upload takes. Seven further candidates were investigated and rejected, two of them
not procurement cases on closer reading. Table~\ref{tab:corpus-summary} summarises the composition; the full list
with citations is in Appendix~\ref{app:corpus}.

\begin{table}[htbp]
\centering
\footnotesize
\caption[Corpus composition]{Composition of the 23-case corpus. Only the merits
group supports an agreement claim; the interim group is reported for completeness
and excluded from every accuracy figure in this chapter.}
\label{tab:corpus-summary}
\begin{tabularx}{\textwidth}{@{}lccX@{}}
\toprule
\textbf{Group} & \textbf{Cases} & \textbf{CA won / lost} & \textbf{Why grouped this way} \\
\midrule
Merits disposition & 14 & 6 / 8 &
 A court ruled on the substantive question this system assesses \\
Interim only & 9 & --- &
 Suspension or summary-judgment rulings; the merits were never decided, so
 comparing a merits prediction to them is not a check \\
\bottomrule
\end{tabularx}
\end{table}

The corpus grew in three stages, and what constrained it changed each time. On a
CPU-bound local model a single negotiation took around seventeen minutes, making a
twenty-case corpus a compute problem. Once inference moved to a tunnelled A10G GPU
that latency fell roughly thirty-fold and the bottleneck became research and
verification: the largest expansion added fifteen cases in one session, the
compute finishing inside an hour and the checking taking the rest.

\section{Structural reliability}
\label{sec:structural}

Across the 50 batches carrying the instrumentation described in
\S\ref{sec:compliance-method}, 2{,}861 of 2{,}865
structured responses required no repair, giving a corpus-wide structural
compliance rate of \textbf{0.9986}. Forty-eight of the 50 batches were clean at
1.00; the worst single batch scored 0.906 and predates the repair-retry pattern
described below.

That pattern is itself a finding. Adding the anti-fabrication and
governing-legislation instructions to the pre-negotiation prompts measurably
raised the rate at which the model dropped a required JSON field, first showing as
three failures in an eight-run batch where the prior corpus had none. The cause
was not the instruction's content but its length: a longer, denser prompt degraded
adherence to the output-format instruction at the end of it, the positional
sensitivity Lu et al.\ \cite{lu2022fantastically} document for few-shot prompts,
appearing here in a format instruction rather than an example ordering. The fix
was a single repair retry re-sending the message with the validation error
appended. The interaction is worth naming for anyone hardening prompts against a
small model: \emph{adding an instruction has a cost paid somewhere else in the
same prompt}, invisible unless something counts it.

\section{Agreement with real dispositions}
\label{sec:agreement}

On the 14 cases with a genuine merits disposition, the simulated outcome
direction matches the real one in 13. That figure needs splitting before it can be
relied on, and \S\ref{sec:leakage} does the splitting: on the eight cases where
the scenario does not state what the real court decided, agreement is
\textbf{7 of 8}, and that is the number this report defends. The exception is Turning Point v Norfolk County Council, and it is not a bug.
Reading the transcript rather than the outcome distribution, the simulated Court's
reasoning is almost verbatim the real claimant's losing argument: that the
authority should have sought clarification of a caveated tender before rejecting
it. The real judge rejected exactly that. Nothing was invented and the reasoning
engages the real facts throughout; the model reached a different legal conclusion
on a genuinely arguable point, a category of error \S\ref{sec:fabrication}
separates from fabrication.

Two caveats attach. On Faraday the direction is correct but the reasoning is not
trustworthy, for reasons given below. On Consultant Connect the direction is
correct but the remedy shape is wrong: the real court imposed financial penalties
on three NHS bodies, which this system's vocabulary cannot express, so
``re-evaluation'' stands in for ``the authority was at fault''. The same gap shows
on Faraday, whose real remedy was a declaration of ineffectiveness, and on Woods,
where the court determined the corrected winner itself rather than merely ordering
the scoring redone. The vocabulary was built for scoring challenges, and it
shows.

Across all 299 runs the outcome distribution comes out heavily skewed
(Figure~\ref{fig:outcomes}). One run produced ``request for additional
information'', a string outside the closed vocabulary. That is a single
instance in 299, from an early four-run batch, and precisely the kind of drift the
downstream module is designed to raise on rather than quietly average away. Median
resolution takes one round; the maximum observed is five.

\begin{figure}[htbp]
\centering
\resizebox{\textwidth}{!}{%
\begin{tikzpicture}[font=\footnotesize]
 % 241 maps to 10cm; every bar is direct-labelled so identity never rests on colour
 \foreach \i/\lab/\cnt/\len/\pct in {
 0/{re-evaluation}/241/10.000/{80.6\%},
 1/{no remedy, decision stands}/31/1.286/{10.4\%},
 2/{deadlock at round limit}/17/0.705/{5.7\%},
 3/{disclosure ordered}/9/0.373/{3.0\%}} {
 \pgfmathsetmacro\y{-\i*0.62}
 \node[anchor=east,font=\scriptsize] at (-0.15,\y) {\lab};
 \fill[f21green] (0,{\y-0.16}) rectangle ({\len},{\y+0.16});
 \node[anchor=west,font=\scriptsize,f21ink] at ({\len+0.12},\y) {\pct\ \ (\cnt)};
 }
 % the out-of-vocabulary run, marked as a different kind of thing
 \node[anchor=east,font=\scriptsize] at (-0.15,-2.48) {request for additional information};
 \fill[f21amber] (0,-2.64) rectangle (0.041,-2.32);
 \node[anchor=west,font=\scriptsize,f21amber] at (0.16,-2.48)
 {0.3\%\ \ (1) --- outside the closed vocabulary};
 \draw[f21grey,line width=0.4pt] (0,-2.80) -- (10.6,-2.80);
 \foreach \x/\t in {0/0, 2.075/50, 4.150/100, 6.224/150, 8.299/200, 10/241} {
 \draw[f21grey,line width=0.4pt] (\x,-2.80) -- (\x,-2.92);
 \node[anchor=north,font=\scriptsize,f21grey] at (\x,-2.94) {\t};
 }
 \node[anchor=north,font=\scriptsize,f21grey] at (5,-3.35) {runs ($n=299$)};
\end{tikzpicture}}
\caption[Outcome distribution]{Outcome distribution across all 299 logged
negotiations. Four in-vocabulary outcomes (green) and the single run that produced
a string outside the Court agent's closed remedy vocabulary (amber). The skew
toward re-evaluation is what makes the majority-class baseline in
\S\ref{sec:baselines} a demanding floor rather than a straw man.}
\label{fig:outcomes}
\end{figure}

\section{Outcome leakage: is the simulation being told the answer?}
\label{sec:leakage}

No pre-processing step strips the real court result out of a source document
before the agents see it, and that bears directly on how the agreement figure
above can be read. The source texts are complete case summaries, and for some
cases extraction carries the disposition through into
\texttt{scenario.description}, the field inserted verbatim into every agent's
prompt. Compounding it, the Court prompt instructs the agent that its recommended
action ``must not end up contradicting an outcome the scenario itself already
tells you occurred''. Where a disposition is present the Court is therefore not
merely informed of the answer but directed towards it.

Auditing all 23 cached scenarios against a deliberately broad pattern
(\texttt{analyze\_outcome\_leakage.py}) finds the disposition stated in eight, six
of them merits cases. Bromcom's description contains ``the court found in favour
of Bromcom on all three grounds''. That is not procedural context, it is the
answer.

\begin{table}[htbp]
\centering
\footnotesize
\caption[Agreement split by outcome leakage]{Directional agreement on merits cases,
split by whether the scenario states the real disposition. The combined figure is
the 13 of 14 quoted above; the split shows how much of it is prediction.}
\label{tab:leakage}
\begin{tabularx}{\textwidth}{@{}>{\raggedright\arraybackslash}p{3.4cm}ccX@{}}
\toprule
\textbf{Subset} & \textbf{Cases} & \textbf{Agreement} & \textbf{What it measures} \\
\midrule
Disposition stated & 6 & 6 / 6 &
 The Court was told the outcome and instructed not to contradict it.
 Compliance, not prediction \\
Disposition absent & 8 & \textbf{7 / 8} &
 Reasoning from facts alone. The single miss is Turning Point \\
Combined & 14 & 13 / 14 & Should not be quoted without the split \\
\bottomrule
\end{tabularx}
\end{table}

The leaked subset scoring 6 of 6 is close to uninformative: an agent told the
answer and instructed to respect it should agree, and does. The leak-free 7 of 8 is
the real result, and a stronger claim than the combined figure precisely because
nothing was given away. The split is not a convenient one either: those eight cases
divide four authority wins to four losses, so the majority-class floor is 4 of 8
and 7 of 8 clears it by a real margin. The leak was measured rather than removed
--- stripping dispositions at ingestion is straightforward for this corpus but not
for arbitrary uploads --- and every accuracy claim rests on the subset where it
does not apply.

\section{Baselines and ablations: does the architecture earn itself?}
\label{sec:baselines}

This is the question an examiner asks first, so it was tested head-on across the
six cases where a merits disposition existed at the time. Reporting a score with
no comparable floor beneath it is the failure Dodge et al.\ \cite{dodge2019showyourwork} single
out, so three baselines were built: a majority-class heuristic with no model call,
a zero-shot prompt deliberately \emph{not} hardened the way the Court prompt is,
and an ablation replacing the Court node with one that cannot resolve.

\begin{table}[htbp]
\centering
\footnotesize
\caption[Baseline comparison]{Directional accuracy against real dispositions, and
the no-Court ablation. The imbalance matters: five of the six cases are ``CA
lost'', so always guessing the majority class scores 5/6 with no reasoning
whatsoever.}
\label{tab:baselines}
\begin{tabularx}{\textwidth}{@{}>{\raggedright\arraybackslash}p{4.3cm}ccX@{}}
\toprule
\textbf{System} & \textbf{Runs} & \textbf{Direction correct} & \textbf{What it establishes} \\
\midrule
Majority-class heuristic & --- & 5 / 6 & The floor. No model call, no reasoning \\
Zero-shot single prompt & 5 per case & 6 / 6 & A naive prompt is not worse here \\
Full pipeline (V4) & 8 per case & 6 / 6 & Ties the naive prompt \\
No-Court ablation & 3 per case & --- & \textbf{0 of 18 runs resolved} \\
\bottomrule
\end{tabularx}
\end{table}

The headline result is negative and reported as measured. \textbf{At $n=6$ the
full pipeline and the naive single-prompt baseline tie: both get every case right,
and the gap down to a predictor with no reasoning is one case wide.} No claim
about the architecture's value can rest on accuracy alone, and establishing that
is the point of the baselines: without them the pipeline's 6/6 would have looked
like design quality, when the majority-class floor shows most of it is the
corpus's class imbalance. A second check applied the Court agent's own grounding
screen to all 30 zero-shot reasoning texts to see whether the unhardened prompt
fabricated more: zero ungrounded numbers. That is a genuine null result, not a
finding bent into a narrative it does not support.

The ablation is a different kind of result. Removing the Court agent produces
deadlock in \textbf{all 18 runs, uniformly on every case} --- not a trend that
could be sampling luck but an architectural necessity, since nothing else in the
graph can decide the parties have converged. An ablation that inferred resolution
from conciliatory dialogue would have defeated the point, which is precisely that
nothing here resolves without adjudication. The defensible claim is therefore
about mechanism, not accuracy: the pipeline produces a modelled negotiation with
declared interests, BATNAs, rounds of concession and a court reviewing the actual
exchange. None of that exists in a single call, and all of it is what the
downstream components depend on.

\section{The V3/V4 prompt ablation}
\label{sec:ablation}

Replacing the Court agent's three-line guiding-principles block with Doc 1's six
judiciary interest categories is a single-variable change, and its effect is
consistent. Every cell is eight repeated runs of one fixed scenario, so the spread
within a cell sizes the sampling noise rather than case
variation. Across every case where V3 deadlocks, V4 reaches zero
deadlock and resolves in roughly half the rounds (Figure~\ref{fig:ablation},
Table~\ref{tab:ablation}).

\begin{figure}[htbp]
\centering
\resizebox{\textwidth}{!}{%
\begin{tikzpicture}[font=\footnotesize]
 % axes
 \draw[f21grey,line width=0.5pt] (0,0) -- (11.6,0);
 \foreach \y/\lab in {0/0, 1.4/.25, 2.8/.50, 4.2/.75, 5.6/1.0} {
 \draw[f21grey!45,line width=0.4pt] (0,\y) -- (11.6,\y);
 \node[anchor=east,f21grey,font=\scriptsize] at (-0.1,\y) {\lab};
 }
 \node[rotate=90,anchor=south,f21grey,font=\scriptsize] at (-0.95,2.8) {resolution rate};

 % five case groups; x centres
 \foreach \i/\lbl/\vthree in {0/Lancashire/4.9, 1/Faraday/4.2, 2/Parkingeye/4.9,
 3/Alstom/2.8, 4/Woods/5.6} {
 \pgfmathsetmacro\xc{0.75 + \i*2.25}
 \fill[f21grey!55] ({\xc-0.62},0) rectangle ({\xc-0.06},\vthree);
 \fill[f21green] ({\xc+0.06},0) rectangle ({\xc+0.62},5.6);
 \node[anchor=north,font=\scriptsize] at (\xc,-0.12) {\lbl};
 }
 % legend
 \fill[f21grey!55] (8.9,6.25) rectangle (9.25,6.55);
 \node[anchor=west,font=\scriptsize] at (9.32,6.4) {V3};
 \fill[f21green] (10.2,6.25) rectangle (10.55,6.55);
 \node[anchor=west,font=\scriptsize] at (10.62,6.4) {V4};
\end{tikzpicture}}
\vspace{4mm}
\caption[V3 vs V4 resolution rates]{Resolution rate by Court prompt version,
$n=8$ per bar, five real cases. V4 sits at the ceiling everywhere; V3 varies by
case. Reaching 1.00 on every case is not automatically the better result --- see
the text.}
\label{fig:ablation}
\end{figure}

\begin{table}[htbp]
\centering
\footnotesize
\caption[Ablation significance]{V3/V4 resolution rates with Fisher's exact tests
\cite{fisher1922interpretation}. No case reaches $p<0.05$. At $n=8$ per arm this
test has low power, so these are ``not distinguishable from noise at this sample
size'', not ``no difference'' \cite{dror2018hitchhiker}.}
\label{tab:ablation}
\begin{tabular}{@{}lcccc@{}}
\toprule
\textbf{Case} & \textbf{V3} & \textbf{V4} & \textbf{V3 avg.\ rounds} & \textbf{Fisher $p$} \\
\midrule
Lancashire Care & 0.875 & 1.000 & 2.75 & 1.000 \\
Faraday & 0.750 & 1.000 & 3.00 & 0.467 \\
Parkingeye & 0.875 & 1.000 & 1.50 & 1.000 \\
Alstom & 0.500 & 1.000 & 4.38 & 0.077 \\
Woods & 1.000 & 1.000 & 1.00 & 1.000 \\
\bottomrule
\end{tabular}
\end{table}

Two things follow from this, and they pull in opposite directions. The pattern
is real and holds consistently across four independent cases, and it matches the
mechanism you'd predict from the change: Doc 1's fifth judicial category is
efficient use of judicial resources, which favours settlement and narrowing of
issues. But nothing here reaches statistical significance, and a prompt that
resolves everything isn't obviously the more correct one for that reason alone.
On Alstom, V4's manifest-error detection rate is identical to V3's at 0.50, so
the difference lies in what happens \emph{after} a finding, not in
whether one gets made. The honest reading is that V4 changes remedy
efficiency rather than judgement, which is what its explicit precedence note was
written to ensure but ``V4 improves resolution rate'' is not a claim that should
travel without that qualification attached.

\section{Fabrication: three distinct failure modes}
\label{sec:fabrication}

This is the project's central concern, and the screens are correspondingly
paranoid. Three independent checks run over the corpus
(Table~\ref{tab:fabrication}); the failure modes they separate are shown in
Figure~\ref{fig:failures}.

\begin{table}[htbp]
\centering
\footnotesize
\caption[Fabrication screens]{The three grounding screens, over 299 logs.}
\label{tab:fabrication}
\begin{tabularx}{\textwidth}{@{}>{\raggedright\arraybackslash}p{3.6cm}cX@{}}
\toprule
\textbf{Screen} & \textbf{Result} & \textbf{Reading} \\
\midrule
Numeric grounding (Step 2B) & 9 flagged of 425 (2.1\%) &
 All nine re-read individually; none is a Court-originated fabrication
 \textbf{within the screened population}. One known miss sits outside it \\
Legal citation, regime & 162 of 702 (23.1\%) wrong regime &
 PCR 2015 or the EU Directive cited in a Procurement Act 2023 world \\
Legal citation, s.12 content & 7 of 21 (33.3\%) misattributed &
 Subsection contents swapped, e.g.\ transparency assigned to s.12(1)(a) \\
\bottomrule
\end{tabularx}
\end{table}

The numeric screen extracts every score-shaped number from a qualitative
assessment and checks whether it appears in the scenario description. Nine of 425
were flagged, and reading all nine changes the interpretation completely: none is
the Court inventing a figure. Six reproduce a pattern first seen on Alstom, where
the negotiating agents invent quantitative specificity and the Court declines to
treat it as verified; two are the Court quoting a number attributively; one
repeats a party's own proposal. The clearest is worth quoting, because the
detector fired on the discipline working as designed:

\begin{quote}\itshape\small
``I do not have enough information to perform a calculation myself, as no specific
numbers or formula are provided in the scenario.''
\hfill --- Court agent, Alstom, round 2
\end{quote}

That clean result holds only inside the population the screen examines, and one
case falls outside it. On Bromcom the Court opened ``this scenario is numeric'',
invented a complete per-bidder dataset (``let's say the scores are as follows:
Bidder 1: A=80, B=70, C=90\ldots''), and performed real arithmetic on twelve values
that appear nowhere in the case. The screen reports nothing, on two independent
grounds: the run classified itself \emph{numeric}, so it never entered the 425
qualitative assessments, and the score-shaped regex matches ``80\%'' or ``scored
80'' but not ``A=80'', so it would have extracted nothing even had the run been
screened. Both are limits of the detector rather than of the Court prompt, and
they bound the finding above as much as they qualify it: the honest claim is that
no Court-originated fabrication survives \emph{this} screen on \emph{these} 425
assessments, not that none occurred.

The screen's real signal sits further upstream: the negotiating agents fabricate
factual premises, and the Court adopts them. The clearest instance is Faraday,
which is not a scoring dispute at all. The authority structured a
\pounds125\,m regeneration deal so as to avoid running a competition, and Faraday
never submitted a bid to be scored. Extraction captured this correctly, labelling
the dispute \texttt{process\_avoidance}; the negotiating agents did not preserve
it. In 3 of 8 V3 runs and 6 of 8 V4 runs the bidder asserts a ``scoring
discrepancy between our bid and that of St Modwen Developments Ltd'', a comparison
that cannot exist since neither party was ever scored. The cause is traceable:
every CA and bidder prompt was developed against scored-bid disputes, and neither
has a branch for a case where no bid was scored. Confronted with an unfamiliar
dispute type, both fall back on the only script they have.

\begin{figure}[htbp]
\centering
\resizebox{\textwidth}{!}{%
\begin{tikzpicture}[font=\footnotesize,
 m/.style={draw,rounded corners=2pt,line width=0.9pt,align=left,inner sep=6pt,
 text width=3.55cm,minimum height=3.0cm,anchor=north west}]
 \node[m,draw=f21green,fill=f21green!8] at (0,0)
 {\textbf{1. Numeric fabrication}\\[2pt]
 \scriptsize Inventing a sub-score or formula the case does not state.\\[3pt]
 \textbf{Status:} addressed within the screened population;
 one known miss outside it (Bromcom).};
 \node[m,draw=f21amber,fill=f21amber!10] at (4.05,0)
 {\textbf{2. Premise fabrication}\\[2pt]
 \scriptsize A negotiating agent inventing a fact about the dispute, which the
 Court then adopts.\\[3pt]
 \textbf{Status:} open. 6 of 8 runs on Faraday.};
 \node[m,draw=f21grey,fill=f21mist] at (8.1,0)
 {\textbf{3. Judgment divergence}\\[2pt]
 \scriptsize Grounded reasoning reaching a different legal conclusion than the
 real judge.\\[3pt]
 \textbf{Status:} not a defect. Turning Point.};
\end{tikzpicture}}
\caption[Failure mode taxonomy]{Three failure modes that a single ``hallucination''
label would conflate. Hardening the adjudicator addresses the first and does
nothing about the second, because the second enters the system before the
adjudicator sees it.}
\label{fig:failures}
\end{figure}

\subsection{Legal citation}
\label{sec:citation}

Citation is the same discipline in a different field, and holds up far less
well. Of 702 citations logged, 162 (23.1\%) name the wrong regime entirely, citing
the Public Contracts Regulations 2015 or the EU Directive inside a system grounded
in the Procurement Act 2023; one names a ``Public Contracts Regulations 2020''
that does not exist. Of the 21 citations specific enough to check against section
12, 7 misattribute the subsection's contents. Both checks were verified against
legislation.gov.uk rather than assumed (Figure~\ref{fig:citations}). The remaining
519 are bucketed as not specific enough to check rather than silently passed as
correct, a deliberately conservative choice: the true error rate is bounded below
by these figures, not estimated by them.

\begin{figure}[htbp]
\centering
\resizebox{\textwidth}{!}{%
\begin{tikzpicture}[font=\footnotesize]
 % 702 citations mapped onto 12cm, 2px surface gap between adjacent segments
 \fill[f21green] (0,0) rectangle (0.239,0.62);
 \fill[f21amber] (0.279,0) rectangle (3.169,0.62);
 \fill[f21grey!55] (3.209,0) rectangle (12.0,0.62);
 \draw[f21ink,line width=0.4pt] (0,0) rectangle (12.0,0.62);

 % direct labels, so the segments are never identified by colour alone
 \draw[f21green,line width=0.7pt] (0.12,0.62) -- (0.12,1.30);
 \node[anchor=south west,font=\scriptsize,f21green,align=left] at (-0.05,1.32)
 {\textbf{14} verified correct\\ against s.12 (2.0\%)};

 \draw[f21amber,line width=0.7pt] (1.72,0.62) -- (1.72,2.30);
 \node[anchor=south,font=\scriptsize,f21amber,align=center] at (1.72,2.32)
 {\textbf{169} checked and wrong (24.1\%)\\
 162 wrong regime + 7 s.12 misattributed};

 \draw[f21grey,line width=0.7pt] (7.6,0) -- (7.6,-0.85);
 \node[anchor=north,font=\scriptsize,f21grey,align=center] at (7.6,-0.87)
 {\textbf{519} not specific enough to check (73.9\%)\\
 bucketed as unknown, deliberately not counted as correct};
\end{tikzpicture}}
\caption[Citation validity]{How the 702 legal citations produced across 299
negotiations divide. Only 183 were specific enough to verify, and of those
169 were wrong. The grey block is the reason the headline 23.1\% wrong-regime
figure is a lower bound rather than an estimate: those citations were not shown to
be correct, only found too general to check.}
\label{fig:citations}
\end{figure}

This is the clearest limitation of prompt-only grounding. The citation-discipline
instruction cut fabricated pinpoint citations, but cannot hand the model knowledge
of which regime governs a 2018 procurement, and an 8B model's priors are saturated
with pre-2023 material whatever the prompt says. Retrieval \cite{lewis2020rag}
over the legislation is the right fix, and the strongest argument in the project
for building the layer that was scoped out.

\section{Negotiation dynamics and BATNA}
\label{sec:dynamics}

Two behavioural checks confirm the negotiation is an actual negotiation rather
than two monologues in parallel. The authority's concession rate rises from 0.013
in round one to 0.470 later; the bidder concedes far less, 0.000 to 0.075,
matching what a bidder concession amounts to: dropping the challenge entirely, a
much larger move than offering extra transparency (Figure~\ref{fig:concession}).
Consecutive same-role message similarity averages 0.097 over 426 pairs, so the
repetition suppression is working.

\begin{figure}[htbp]
\centering
\resizebox{0.82\textwidth}{!}{%
\begin{tikzpicture}[font=\footnotesize]
 % y: concession rate 0 to 0.5 over 5cm. x: round 1 at 0, round 2+ at 7
 \foreach \v/\ylen in {0/0, 0.1/1, 0.2/2, 0.3/3, 0.4/4, 0.5/5} {
 \draw[f21grey!35,line width=0.4pt] (0,\ylen) -- (7,\ylen);
 \node[anchor=east,font=\scriptsize,f21grey] at (-0.12,\ylen) {\v};
 }
 \node[anchor=south,font=\scriptsize] at (0,5.45) {round 1};
 \node[anchor=south,font=\scriptsize] at (7,5.45) {rounds 2+};
 \node[rotate=90,anchor=south,font=\scriptsize,f21grey] at (-0.95,2.5)
 {concession rate};

 % Contracting Authority: 0.013 -> 0.470
 \draw[f21green,line width=2pt] (0,0.13) -- (7,4.70);
 \fill[f21green] (0,0.13) circle (3.4pt);
 \fill[f21green] (7,4.70) circle (3.4pt);
 \node[anchor=west,font=\scriptsize,f21green] at (7.2,4.70)
 {\textbf{Contracting Authority} 0.470};
 \node[anchor=east,font=\scriptsize,f21green] at (-0.35,0.13) {0.013};

 % Aggrieved Bidder: 0.000 -> 0.075
 \draw[f21blue,line width=2pt] (0,0.0) -- (7,0.75);
 \fill[f21blue] (0,0.0) circle (3.4pt);
 \fill[f21blue] (7,0.75) circle (3.4pt);
 \node[anchor=west,font=\scriptsize,f21blue] at (7.2,0.75)
 {\textbf{Aggrieved Bidder} 0.075};
 \node[anchor=east,font=\scriptsize,f21blue] at (-0.35,-0.30) {0.000};
\end{tikzpicture}}
\caption[Concession dynamics]{Concession rate by role, round 1 against rounds 2
and later, over 299 negotiations. The authority opens almost immovably and then
concedes about half the time; the bidder barely moves at either stage. The
asymmetry is not a modelling artefact but a consequence of what a concession
means on each side, since the only real concession available to a bidder is
dropping its challenge.}
\label{fig:concession}
\end{figure}

Zero explicit
retractions were found, using a check that was narrowed after an earlier version
misread ``willing to withdraw our claim'', which is a settlement offer, as a
retraction.

Table~\ref{tab:batna} shows how each agent judges its own outcome against its
declared BATNA. The asymmetry here is structural and matches the law: an
authority that successfully defends its award loses nothing, while a bidder that
obtains a re-evaluation has still spent the bid cost and may well lose again.

\begin{table}[htbp]
\centering
\footnotesize
\caption[BATNA self-assessment]{Self-reported outcome relative to declared BATNA,
$n=299$ per role, classified by a documented keyword heuristic. The heuristic is
not a ground-truth label, and the large ``unclear'' fractions are reported rather
than redistributed.}
\label{tab:batna}
\begin{tabular}{@{}lcccc@{}}
\toprule
\textbf{Role} & \textbf{Beats} & \textbf{Matches} & \textbf{Falls short} & \textbf{Unclear} \\
\midrule
Contracting Authority & 71\% & 1\% & 2\% & 27\% \\
Aggrieved Bidder & 24\% & 0\% & 21\% & 56\% \\
\bottomrule
\end{tabular}
\end{table}

An earlier run over 63 logs reported the authority as never falling short; at
299 it falls short in five cases, all genuine self-reported shortfalls on
inspection. The absolute claim was never safe from the smaller sample, and the
correction is a reminder of how cautiously any proportion here should be cited
before $N$ grows.

\section{Readability: the summariser does not meet its own claim}
\label{sec:readability}

The Summary agent's schema describes its output as ``a 3--4 sentence summary a
non-lawyer could understand''. Over all 299 summaries, mean Flesch Reading Ease
\cite{flesch1948readability} is \textbf{24.7} (median 25.9), the ``very difficult,
graduate level'' band; mean Flesch--Kincaid grade \cite{kincaid1975derivation} is
14.7. The best summary scores 57.5, short of the 60--70 plain-English band, and
the worst is negative. The figure is unchanged from an earlier run over a smaller
corpus, so this is a stable property of the prompt rather than a sampling
artefact. The agent that exists to make the output accessible does not, by a
standard measure, succeed --- and it would have gone unnoticed without measuring,
because the summaries read fluently on the page.

\section{The risk screen against real cases}
\label{sec:risk-eval}

The screen was run against eight real cases, filling a profile field only where
the judgment states or entails it. Six produced any populated field; Faraday and
Alstom produced none, which is the finding this section turns on. Two of the six
reach ``medium'' from bidder-side facts alone: Braceurself, with a stated 80.25\%
against 82.5\% margin, and Lancashire Care, where the court found the authority's
reasons ``too vague and generic''. Both are among the clearest cases where a court
found the process defective. That co-occurrence is worth naming but is not a
validated correlation, and six cases could not establish one.

The remaining four land ``low'' despite being real litigated challenges, which is
a finding about the data rather than a failure of the screen. Judgments record
legal reasoning; they rarely record whether the evaluation panel was
procurement-trained or the audit trail complete, which are exactly the conditions
the rules key on. The module already states that an unpopulated profile always
scores low, and that this is a gap in the input rather than a finding.

\section{Threats to validity}
\label{sec:threats}

\textbf{Sample size.} Every rate reported here is a proportion over a modest
$N$. The ablation is $n=8$ per cell with no significant result; the baseline
comparison is $n=6$. The deadlock finding is the exception, holding at 18 of 18
by architectural necessity rather than sampling luck.

\textbf{Outcome leakage.} Six of the 14 merits scenarios state the real
disposition in the text the agents receive, and the Court prompt instructs the
agent not to contradict it (\S\ref{sec:leakage}). No accuracy claim rests on those
six.

\textbf{Single model, single configuration.} Everything runs on Llama 3.1 8B at
fixed per-role temperatures. Nothing here separates ``this architecture'' from
``this model under these particular prompts''.

\textbf{Corpus bias.} Litigated cases are already a filtered minority, and this
corpus is filtered further toward cases documented well enough to verify. The
merits group splits 6 authority wins to 8 losses, but that balance was achieved by
deliberately hunting for authority-win cases, which is itself a selection
choice.

\textbf{Detector limits.} The numeric screen is string-presence matching against
the scenario description, not a fact-checker, and is known to have missed one real
Court fabrication on both of its axes (\S\ref{sec:fabrication}); the Faraday
premise-fabrication count comes from a phrase search and is a lower bound; the
BATNA classifier is a keyword heuristic. Each is stated plainly in its own module.

\textbf{Structural mismatch.} Nine of 23 cases were decided at an interim stage
this system cannot model. Any apparent correspondence on those cases is
directional at best, and they were excluded from every accuracy figure in this
chapter.

\section{Reflection on the process}
\label{sec:reflection}

Three things went differently from the plan, each productively.

First, the interesting problem moved. The plan treated the negotiation system as
the deliverable and evaluation as a final phase; in practice the negotiation
worked within weeks and the remaining months went into evaluation, because the
question that mattered was whether anything it produced could be trusted.

Second, the client reframed the project in a sentence. Fusion21's interest was
prevention, not replay, and the risk screen exists because of it: the smallest
component, no model call, and the only part with an observable ground truth in
principle. That makes it the most scientifically tractable thing here, and the
thing this project is least able to validate without data only the client can
supply.

The third is that the most useful bugs came from documents I had not chosen. A
board demonstration ended with a Fusion21 employee handing me an award letter and
a scoring table and asking what the system made of them. It found a genuine
arithmetic defect: the Court re-applying an already-applied percentage weighting
to a sub-score, turning a correct evaluation into a false manifest-error finding
that three increasingly explicit prompt revisions failed to fix. The model would
state the correct approach in words, then execute the wrong arithmetic in the same
response. Those revisions were reverted rather than shipped, and the limitation is
recorded as open. Choosing my own test inputs had been selecting for the failures
I already knew about.